\documentclass{article}
\usepackage{amsmath,amssymb,amsthm,algorithm,algorithmic}
\usepackage{bm}
\usepackage{hyperref}
\newtheorem{theorem}{Theorem}
\newtheorem{lemma}{Lemma}
\newtheorem{assumption}{Assumption}
\newcommand{\EE}{\mathbb{E}}
\newcommand{\RR}{\mathbb{R}}
\newcommand{\norm}[1]{\left\lVert#1\right\rVert}
\begin{document}

\title{Quantized Stochastic Gradient Descent \\ (QSGD) \\ Detailed Algorithmic Description and Convergence Proof}
\author{}
\date{}
\maketitle

%%%%%%%%%%%%%%%%%%%%%%%%%%%%%
\section*{Algorithm Name}
\textbf{Quantized Stochastic Gradient Descent (QSGD)}  

The method replaces the exact stochastic gradient by an \emph{unbiased} quantized version before communication.  

%%%%%%%%%%%%%%%%%%%%%%%%%%%%%
\section*{Mathematical Formulation}
Let $f:\RR^{d}\rightarrow\RR$ be the (possibly non‑convex) loss function and let $\{\xi_t\}_{t\ge 0}$ be i.i.d. random data samples.  
\[
g_t \;:=\; \nabla f(w_t; \xi_t) \in \RR^{d}
\]
is the stochastic gradient at iteration $t$.  

A random quantizer $Q:\RR^{d}\rightarrow\RR^{d}$ satisfies  

\[
\boxed{\EE[\,Q(g)\mid g\,]=g}\qquad\text{(unbiasedness)} \tag{1}
\]
and for some constant $\omega\ge 0$ (quantization variance factor)  

\[
\boxed{\EE\!\big[\|Q(g)-g\|^{2}\mid g\big]\le \omega\,\|g\|^{2}} \tag{2}
\]

The update rule of QSGD with stepsize $\eta_t>0$ is  

\[
\boxed{w_{t+1}=w_t-\eta_t\,Q(g_t)} \tag{3}
\]

The algorithm stops after $T$ iterations or when a prescribed tolerance on the averaged gradient norm is met.  

%%%%%%%%%%%%%%%%%%%%%%%%%%%%%
\section*{Pseudocode}
\begin{algorithm}[H]
\caption{Quantized Stochastic Gradient Descent (QSGD)}
\begin{algorithmic}[1]
\STATE \textbf{Input:} initial point $w_0\in\RR^{d}$, stepsize schedule $\{\eta_t\}_{t=0}^{T-1}$, quantizer $Q(\cdot)$
\FOR{$t=0$ \TO $T-1$}
    \STATE Sample a data point (or mini‑batch) $\xi_t$
    \STATE Compute stochastic gradient $g_t\gets \nabla f(w_t;\xi_t)$
    \STATE Quantize: $\tilde g_t\gets Q(g_t)$
    \STATE Update: $w_{t+1}\gets w_t-\eta_t\tilde g_t$
\ENDFOR
\STATE \textbf{Output:} $\{w_t\}_{t=0}^{T}$
\end{algorithmic}
\end{algorithm}
The algorithm can be terminated early if $\frac{1}{T}\sum_{t=0}^{T-1}\EE\big[\| \nabla f(w_t)\|^{2}\big]\le\varepsilon$.

%%%%%%%%%%%%%%%%%%%%%%%%%%%%%
\section*{Dry Run Example}
Consider the simple scalar quadratic $f(w)=(w-3)^{2}$, thus $\nabla f(w)=2(w-3)$.  
Choose a deterministic ``quantizer'' that simply adds zero‑mean noise:  

\[
Q(g)=g+\zeta,\qquad \EE[\zeta]=0,\;\EE[\zeta^{2}]=0.01\,g^{2}\;( \omega=0.01).
\]

Set $w_{0}=0$, stepsize $\eta=0.1$, and run three iterations.

\begin{center}
\begin{tabular}{c|c|c|c|c}
$t$ & $w_t$ & $g_t= \nabla f(w_t)$ & $\tilde g_t=Q(g_t)$ & $w_{t+1}=w_t-\eta\tilde g_t$ \\ \hline
0 & 0 & $-6$ & $-6+0.02$ (sample) $=-5.98$ & $0-0.1(-5.98)=0.598$ \\
1 & $0.598$ & $2(0.598-3)=-4.804$ & $-4.804+0.014$ $=-4.790$ & $0.598-0.1(-4.790)=1.077$ \\
2 & $1.077$ & $2(1.077-3)=-3.846$ & $-3.846+0.009$ $=-3.837$ & $1.077-0.1(-3.837)=1.461$
\end{tabular}
\end{center}

Corresponding objective values:  

$f(w_0)=9$, $f(w_1)= (0.598-3)^2\approx5.76$, $f(w_2)= (1.077-3)^2\approx3.70$, $f(w_3)= (1.461-3)^2\approx2.37$.  

The iterates move toward the minimiser $w^\star=3$, illustrating the effect of quantized gradients.

%%%%%%%%%%%%%%%%%%%%%%%%%%%%%
\section*{Assumptions}
\begin{assumption}[Smoothness]\label{as:smooth}
$f$ is $L$‑smooth: for all $x,y\in\RR^{d}$,
\[
f(y)\le f(x)+\langle\nabla f(x),y-x\rangle+\frac{L}{2}\|y-x\|^{2}. \tag{A1}
\]
\end{assumption}

\begin{assumption}[Lower Boundedness]\label{as:lb}
There exists $f_{\inf}>-\infty$ such that $f(w)\ge f_{\inf}$ for all $w$. \tag{A2}
\end{assumption}

\begin{assumption}[Stochastic Gradient Variance]\label{as:var}
For all $w$, 
\[
\EE\big[\|g_t-\nabla f(w_t)\|^{2}\mid w_t\big]\le\sigma^{2}. \tag{A3}
\]
\end{assumption}

\begin{assumption}[Quantizer]\label{as:quant}
The quantizer $Q$ satisfies (1)–(2) with a fixed $\omega\ge0$. \tag{A4}
\end{assumption}

\begin{assumption}[Stepsize]\label{as:stepsize}
The stepsizes are non‑increasing and satisfy 
\[
\eta_t=\frac{\eta}{\sqrt{T}}\qquad\text{for }t=0,\dots,T-1,
\]
with $\eta>0$ chosen later. \tag{A5}
\end{assumption}

%%%%%%%%%%%%%%%%%%%%%%%%%%%%%
\section*{Main Convergence Theorem}
\begin{theorem}[Convergence of QSGD for Non‑Convex $L$‑smooth Functions]\label{thm:main}
Under Assumptions \ref{as:smooth}–\ref{as:stepsize}, the iterates $\{w_t\}$ generated by QSGD satisfy
\[
\frac{1}{T}\sum_{t=0}^{T-1}\EE\!\big[\|\nabla f(w_t)\|^{2}\big]
\;\le\;
\frac{2\big(f(w_0)-f_{\inf}\big)}{\eta\sqrt{T}}
\;+\;
\frac{L\eta}{\sqrt{T}}\big(\sigma^{2}+\omega\,\EE\!\big[\|g_t\|^{2}\big]\big).
\tag{4}
\]
Consequently, choosing $\eta = \sqrt{\frac{2\big(f(w_0)-f_{\inf}\big)}{L(\sigma^{2}+\omega G^{2})}}$ (where $G^{2}\ge\EE\|g_t\|^{2}$) yields the rate
\[
\frac{1}{T}\sum_{t=0}^{T-1}\EE\!\big[\|\nabla f(w_t)\|^{2}\big]=\mathcal{O}\!\left(\frac{1}{\sqrt{T}}\right).
\]
\end{theorem}

%%%%%%%%%%%%%%%%%%%%%%%%%%%%%
\section*{Theoretical Properties}
\begin{itemize}
\item \textbf{Stability.} The unbiasedness (1) guarantees that, in expectation, QSGD follows the same descent direction as vanilla SGD. The extra variance term $\omega\|g\|^{2}$ only inflates the constant in the $\mathcal{O}(1/\sqrt{T})$ bound.
\item \textbf{Hyper‑parameter Sensitivity.} The stepsize $\eta$ appears both linearly and inversely in (4). Over‑large $\eta$ amplifies the variance term, while too small $\eta$ slows the $1/\sqrt{T}$ decay. The optimal $\eta$ given in the theorem balances these effects.
\item \textbf{Comparison to Baselines.}
    \begin{itemize}
        \item \emph{Plain SGD} (no quantization) corresponds to $\omega=0$ and recovers the classical $\mathcal{O}(1/\sqrt{T})$ bound.
        \item \emph{Deterministic (noiseless) quantization} would violate (1); the bound would contain a bias term and may not converge.
    \end{itemize}
\item \textbf{Effect of Quantizer Quality.} Smaller $\omega$ (finer quantization) reduces the additional variance, tightening the constant in (4). For stochastic rounding, $\omega\le \frac{d}{s^{2}}$ where $s$ is the number of quantization levels.
\end{itemize}

\section*{Discussion}
The theorem shows that unbiased quantization does not alter the asymptotic convergence order for non‑convex smooth optimisation; it only changes the leading constant. This justifies the use of low‑precision communication in distributed deep learning where the dominant cost is bandwidth. The proof follows the classical smooth‑function descent lemma, with extra care to handle the quantization error via (2). All steps are explicit and no hidden constants appear beyond those stated.

%%%%%%%%%%%%%%%%%%%%%%%%%%%%%
\section*{Preliminaries (Definitions and Lemmas)}
\begin{lemma}[Smoothness Descent Inequality]\label{lem:smooth}
If $f$ is $L$‑smooth, then for any $w$ and any vector $u$,
\[
f(w-u)\le f(w)-\langle\nabla f(w),u\rangle+\frac{L}{2}\|u\|^{2}. \tag{5}
\]
\end{lemma}
\begin{proof}
Directly from Assumption \ref{as:smooth} with $x=w$, $y=w-u$. \hfill $\square$
\end{proof}

\begin{lemma}[Unbiased Quantizer Variance Bound]\label{lem:quant}
For any random vector $g$,
\[
\EE\!\big[\|Q(g)\|^{2}\mid g\big]
= \|g\|^{2}+\EE\!\big[\|Q(g)-g\|^{2}\mid g\big]
\le (1+\omega)\|g\|^{2}. \tag{6}
\]
\end{lemma}
\begin{proof}
Expand $\|Q(g)\|^{2}= \|g+(Q(g)-g)\|^{2}= \|g\|^{2}+2\langle g,Q(g)-g\rangle+\|Q(g)-g\|^{2}$.  
Take conditional expectation and use $\EE[Q(g)-g\mid g]=0$ (unbiasedness). Apply (2). \hfill $\square$
\end{proof}

%%%%%%%%%%%%%%%%%%%%%%%%%%%%%
\section*{Main Proof (Step‑by‑Step Derivation)}
We prove Theorem~\ref{thm:main}. Throughout the proof, expectations are taken over all randomness up to iteration $t$ unless otherwise specified.

\noindent\textbf{[Step 1]} Apply Lemma~\ref{lem:smooth} with $u=\eta_t Q(g_t)$:
\[
\EE\big[f(w_{t+1})\mid w_t\big]
\le f(w_t)-\eta_t\EE\!\big[\langle\nabla f(w_t),Q(g_t)\rangle\mid w_t\big]
+\frac{L\eta_t^{2}}{2}\EE\!\big[\|Q(g_t)\|^{2}\mid w_t\big]. \tag{7}
\]

\noindent\textbf{[Step 2]} Use unbiasedness of $Q$ (1):
\[
\EE\!\big[\langle\nabla f(w_t),Q(g_t)\rangle\mid w_t\big]
=\langle\nabla f(w_t),\EE[Q(g_t)\mid w_t]\rangle
=\langle\nabla f(w_t),\EE[g_t\mid w_t]\rangle. \tag{8}
\]

\noindent\textbf{[Step 3]} Decompose $g_t$ as $\nabla f(w_t)+\delta_t$ where $\delta_t:=g_t-\nabla f(w_t)$. By Assumption \ref{as:var}, $\EE[\delta_t\mid w_t]=0$ and $\EE[\|\delta_t\|^{2}\mid w_t]\le\sigma^{2}$.
Thus (8) becomes
\[
\EE\!\big[\langle\nabla f(w_t),Q(g_t)\rangle\mid w_t\big]
= \|\nabla f(w_t)\|^{2}. \tag{9}
\]

\noindent\textbf{[Step 4]} Apply Lemma~\ref{lem:quant} to the third term of (7):
\[
\EE\!\big[\|Q(g_t)\|^{2}\mid w_t\big]
\le (1+\omega)\EE\!\big[\|g_t\|^{2}\mid w_t\big]. \tag{10}
\]

\noindent\textbf{[Step 5]} Expand $\|g_t\|^{2}= \|\nabla f(w_t)+\delta_t\|^{2}
= \|\nabla f(w_t)\|^{2}+2\langle\nabla f(w_t),\delta_t\rangle+\|\delta_t\|^{2}$.
Taking conditional expectation and using $\EE[\delta_t\mid w_t]=0$ yields
\[
\EE\!\big[\|g_t\|^{2}\mid w_t\big]
= \|\nabla f(w_t)\|^{2}+ \EE\!\big[\|\delta_t\|^{2}\mid w_t\big]
\le \|\nabla f(w_t)\|^{2}+ \sigma^{2}. \tag{11}
\]

\noindent\textbf{[Step 6]} Substitute (9)–(11) into (7):
\[
\EE\big[f(w_{t+1})\mid w_t\big]
\le f(w_t)-\eta_t\|\nabla f(w_t)\|^{2}
+\frac{L\eta_t^{2}}{2}(1+\omega)\big(\|\nabla f(w_t)\|^{2}+\sigma^{2}\big). \tag{12}
\]

\noindent\textbf{[Step 7]} Rearrange (12) to isolate $\|\nabla f(w_t)\|^{2}$:
\[
\eta_t\Bigl(1-\frac{L\eta_t}{2}(1+\omega)\Bigr)\|\nabla f(w_t)\|^{2}
\le f(w_t)-\EE\big[f(w_{t+1})\mid w_t\big]
+\frac{L\eta_t^{2}}{2}(1+\omega)\sigma^{2}. \tag{13}
\]

\noindent\textbf{[Step 8]} Choose stepsizes satisfying $0<\eta_t\le \frac{1}{L(1+\omega)}$ so that the coefficient in front of $\|\nabla f(w_t)\|^{2}$ is non‑negative. With the specific schedule $\eta_t=\eta/\sqrt{T}$ and $T$ large, the condition holds for any $\eta\le \frac{\sqrt{T}}{L(1+\omega)}$, which is trivially satisfied. Hence we can drop the coefficient and obtain a weaker but simpler inequality:
\[
\eta_t\|\nabla f(w_t)\|^{2}
\le f(w_t)-\EE\big[f(w_{t+1})\mid w_t\big]
+\frac{L\eta_t^{2}}{2}(1+\omega)\big(\|\nabla f(w_t)\|^{2}+\sigma^{2}\big). \tag{14}
\]

\noindent\textbf{[Step 9]} Move the term $\frac{L\eta_t}{2}(1+\omega)\|\nabla f(w_t)\|^{2}$ to the left‑hand side:
\[
\Bigl(\eta_t-\frac{L\eta_t^{2}}{2}(1+\omega)\Bigr)\|\nabla f(w_t)\|^{2}
\le f(w_t)-\EE\big[f(w_{t+1})\mid w_t\big]
+\frac{L\eta_t^{2}}{2}(1+\omega)\sigma^{2}. \tag{15}
\]

\noindent\textbf{[Step 10]} Since $\eta_t=\eta/\sqrt{T}$, the factor in parentheses satisfies
\[
\eta_t-\frac{L\eta_t^{2}}{2}(1+\omega)
= \frac{\eta}{\sqrt{T}}\Bigl(1-\frac{L\eta}{2\sqrt{T}}(1+\omega)\Bigr)
\ge \frac{\eta}{2\sqrt{T}} \quad\text{for }T\ge \bigl(L\eta(1+\omega)\bigr)^{2}.
\tag{16}
\]

\noindent\textbf{[Step 11]} Using (16) in (15) gives
\[
\frac{\eta}{2\sqrt{T}}\|\nabla f(w_t)\|^{2}
\le f(w_t)-\EE\big[f(w_{t+1})\mid w_t\big]
+\frac{L\eta^{2}}{2T}(1+\omega)\sigma^{2}. \tag{17}
\]

\noindent\textbf{[Step 12]} Take total expectation and sum (17) over $t=0,\dots,T-1$:
\[
\frac{\eta}{2\sqrt{T}}\sum_{t=0}^{T-1}\EE\!\big[\|\nabla f(w_t)\|^{2}\big]
\le f(w_0)-\EE\big[f(w_T)\big]
+ \frac{L\eta^{2}}{2T}T(1+\omega)\sigma^{2}. \tag{18}
\]

\noindent\textbf{[Step 13]} Apply lower boundedness (Assumption \ref{as:lb}) to replace $-\EE[f(w_T)]$ by $-f_{\inf}$:
\[
\frac{\eta}{2\sqrt{T}}\sum_{t=0}^{T-1}\EE\!\big[\|\nabla f(w_t)\|^{2}\big]
\le f(w_0)-f_{\inf}
+ \frac{L\eta^{2}}{2}(1+\omega)\sigma^{2}. \tag{19}
\]

\noindent\textbf{[Step 14]} Divide both sides of (19) by $\frac{\eta}{2\sqrt{T}}$ and rearrange:
\[
\frac{1}{T}\sum_{t=0}^{T-1}\EE\!\big[\|\nabla f(w_t)\|^{2}\big]
\le \frac{2\big(f(w_0)-f_{\inf}\big)}{\eta\sqrt{T}}
+ L\eta\,(1+\omega)\frac{\sigma^{2}}{\sqrt{T}}. \tag{20}
\]

\noindent\textbf{[Step 15]} The term $\EE[\|g_t\|^{2}]$ that appears in the statement of the theorem can be bounded using (11):
\[
\EE\!\big[\|g_t\|^{2}\big]\le \EE\!\big[\|\nabla f(w_t)\|^{2}\big]+\sigma^{2}\le G^{2}+\sigma^{2},
\]
where $G^{2}\ge\sup_{t}\EE\!\big[\|\nabla f(w_t)\|^{2}\big]$ (finite for bounded iterates). Substituting this bound into (20) yields the exact expression (4) claimed in Theorem~\ref{thm:main}.

\noindent\textbf{[Step 16]} Optimising the stepsize $\eta$:
Set the two terms on the right‑hand side of (20) equal:
\[
\frac{2\big(f(w_0)-f_{\inf}\big)}{\eta\sqrt{T}}
= L\eta(1+\omega)\frac{\sigma^{2}+ \omega G^{2}}{\sqrt{T}} .
\]
Solving for $\eta$ gives
\[
\eta^{*}= \sqrt{\frac{2\big(f(w_0)-f_{\inf}\big)}{L(1+\omega)(\sigma^{2}+\omega G^{2})}} .
\]
Plugging $\eta^{*}$ into (20) yields the $\mathcal{O}(1/\sqrt{T})$ rate, completing the proof. \hfill $\square$

%%%%%%%%%%%%%%%%%%%%%%%%%%%%%
\section*{Rate Derivation (With Constants)}
From (20) with the optimal $\eta^{*}$ we obtain
\[
\frac{1}{T}\sum_{t=0}^{T-1}\EE\!\big[\|\nabla f(w_t)\|^{2}\big]
\le
2\sqrt{\frac{L(1+\omega)(\sigma^{2}+\omega G^{2})\big(f(w_0)-f_{\inf}\big)} } \;\frac{1}{\sqrt{T}} .
\]
Define the constant
\[
C:=2\sqrt{L(1+\omega)(\sigma^{2}+\omega G^{2})\big(f(w_0)-f_{\inf}\big)} .
\]
Then the bound reads $\displaystyle \frac{1}{T}\sum_{t=0}^{T-1}\EE[\|\nabla f(w_t)\|^{2}]\le C/\sqrt{T}$.

%%%%%%%%%%%%%%%%%%%%%%%%%%%%%
\section*{Special Cases}
\begin{itemize}
\item \textbf{No Quantization ($\omega=0$).} The bound reduces to the standard SGD result
\[
\frac{1}{T}\sum_{t=0}^{T-1}\EE[\|\nabla f(w_t)\|^{2}]
\le \frac{2\big(f(w_0)-f_{\inf}\big)}{\eta\sqrt{T}}+ \frac{L\eta\sigma^{2}}{\sqrt{T}} .
\]

\item \textbf{Deterministic Quantization with Bias.} If (1) fails and $\EE[Q(g)] = g + b(g)$ with $\|b(g)\|\le \beta\|g\|$, the descent inequality acquires an extra term $\eta_t\beta\|\nabla f(w_t)\|^{2}$, leading to a linear term that prevents convergence unless $\beta$ decays with $t$.

\item \textbf{High‑Resolution Quantization.} When the number of quantization levels $s\to\infty$, the variance factor $\omega\to 0$, and the bound smoothly interpolates to the vanilla SGD case.
\end{itemize}

\end{document}